\newcommand{\EduWFFiveDiagram}{%
\begin{tikzpicture}[edu activity]
  \node (start) [edu startstop] {Start};
  \node (open) [edu process, below=of start] {Owner opens the organization member-management screen.};
  \node (revokeaction) [edu decision, below=of open] {Revoke a pending invite?};
  \node (invite) [edu process, below=16mm of revokeaction] {Enter an email address and choose the organization role teacher or student.};
  \node (callinvite) [edu process, below=of invite] {Call the organization-member invitation action.};
  \node (existing) [edu decision, below=of callinvite] {Existing profile already uses the target email?};
  \node (activate) [edu process, below=of existing] {Activate the organization membership immediately.};
  \node (pending) [edu decision, right=42mm of existing] {Pending invite already exists for the same email?};
  \node (refreshinvite) [edu process, right=42mm of pending] {Refresh the existing token and expiry for the pending invite.};
  \node (newinvite) [edu process, below=of pending] {Create a new pending invite token for the organization.};
  \node (revoke) [edu process, right=42mm of revokeaction] {Select a pending invite and call the invite-revocation action.};
  \node (stillpending) [edu decision, below=of revoke] {Invite still pending?};
  \node (revoked) [edu process, below=of stillpending] {Revoke the pending invite and remove it from the active invite list.};
  \node (refresh) [edu process, below=22mm of activate] {Refresh the user list and the pending invites list.};
  \node (error) [edu process, right=42mm of stillpending] {Show the revocation error and leave the invite unchanged.};
  \node (end) [edu startstop, below=of refresh] {End};

  \path [edu arrow] (start) -- (open);
  \path [edu arrow] (open) -- (revokeaction);
  \path [edu arrow] (revokeaction) -- node[edu label] {No} (invite);
  \path [edu arrow] (revokeaction.east) -- node[edu label] {Yes} (revoke.west);
  \path [edu arrow] (invite) -- (callinvite);
  \path [edu arrow] (callinvite) -- (existing);
  \path [edu arrow] (existing) -- node[edu label] {Yes} (activate);
  \path [edu arrow] (existing.east) -- node[edu label] {No} (pending.west);
  \path [edu arrow] (pending) -- node[edu label] {No} (newinvite);
  \path [edu arrow] (pending.east) -- node[edu label] {Yes} (refreshinvite.west);
  \path [edu arrow] (activate) -- (refresh);
  \path [edu arrow] (newinvite.south) |- (refresh.east);
  \path [edu arrow] (refreshinvite.south) |- (refresh.east);
  \path [edu arrow] (revoke) -- (stillpending);
  \path [edu arrow] (stillpending) -- node[edu label] {Yes} (revoked);
  \path [edu arrow] (stillpending.east) -- node[edu label] {No} (error.west);
  \path [edu arrow] (revoked.south) |- (refresh.east);
  \path [edu arrow] (refresh) -- (end);
  \path [edu arrow] (error.south) |- (end.east);
\end{tikzpicture}%
}

\newcommand{\EduWFSixDiagram}{%
\begin{tikzpicture}[edu activity]
  \node (start) [edu startstop] {Start};
  \node (open) [edu process, below=of start] {Owner opens the features and extensions screen.};
  \node (load) [edu process, below=of open] {Load the feature definitions together with the current organization settings.};
  \node (builtin) [edu decision, below=of load] {Configure a built-in feature?};
  \node (toggle) [edu process, below=16mm of builtin] {Toggle the built-in organization feature on or off.};
  \node (parent) [edu decision, below=of toggle] {Parent feature enabled?};
  \node (savefeature) [edu process, below=of parent] {Upsert the built-in feature state to the organization settings table.};
  \node (locked) [edu process, right=42mm of parent] {Keep the child feature locked until the required parent feature is enabled.};
  \node (extension) [edu process, right=42mm of builtin] {Enter the custom extension name, URL, and description.};
  \node (extvalid) [edu decision, below=of extension] {Extension fields valid enough to save?};
  \node (saveext) [edu process, below=of extvalid] {Insert the custom extension record and set its enabled state.};
  \node (exterror) [edu process, right=42mm of extvalid] {Show the validation error and keep the extension unchanged.};
  \node (refresh) [edu process, below=22mm of savefeature] {Refresh the organization feature and extension view.};
  \node (end) [edu startstop, below=of refresh] {End};

  \path [edu arrow] (start) -- (open);
  \path [edu arrow] (open) -- (load);
  \path [edu arrow] (load) -- (builtin);
  \path [edu arrow] (builtin) -- node[edu label] {Yes} (toggle);
  \path [edu arrow] (builtin.east) -- node[edu label] {No} (extension.west);
  \path [edu arrow] (toggle) -- (parent);
  \path [edu arrow] (parent) -- node[edu label] {Yes} (savefeature);
  \path [edu arrow] (parent.east) -- node[edu label] {No} (locked.west);
  \path [edu arrow] (extension) -- (extvalid);
  \path [edu arrow] (extvalid) -- node[edu label] {Yes} (saveext);
  \path [edu arrow] (extvalid.east) -- node[edu label] {No} (exterror.west);
  \path [edu arrow] (savefeature) -- (refresh);
  \path [edu arrow] (locked.south) |- (refresh.east);
  \path [edu arrow] (saveext.south) |- (refresh.east);
  \path [edu arrow] (refresh) -- (end);
  \path [edu arrow] (exterror.south) |- (end.east);
\end{tikzpicture}%
}

\newcommand{\EduWFSevenDiagram}{%
\begin{tikzpicture}[edu activity]
  \node (start) [edu startstop] {Start};
  \node (open) [edu process, below=of start] {Owner opens the class-management screen.};
  \node (deleteaction) [edu decision, below=of open] {Delete an existing class?};
  \node (edit) [edu process, below=16mm of deleteaction] {Create a class or open an existing class for editing.};
  \node (submit) [edu process, below=of edit] {Submit class metadata together with the teacher email.};
  \node (valid) [edu decision, below=of submit] {Teacher profile exists and request validates successfully?};
  \node (saveclass) [edu process, below=of valid] {Create or update the class record, preserving memberships when the teacher does not change.};
  \node (settings) [edu process, below=of saveclass] {Upsert the class feature settings and class extension settings.};
  \node (delete) [edu process, right=42mm of deleteaction] {Confirm deletion and call the class-deletion action.};
  \node (removed) [edu process, below=of delete] {Remove the class from the organization workflow.};
  \node (error) [edu process, right=42mm of valid] {Return the validation error and keep the class unchanged.};
  \node (refresh) [edu process, below=22mm of settings] {Refresh the class-management state for the owner.};
  \node (end) [edu startstop, below=of refresh] {End};

  \path [edu arrow] (start) -- (open);
  \path [edu arrow] (open) -- (deleteaction);
  \path [edu arrow] (deleteaction) -- node[edu label] {No} (edit);
  \path [edu arrow] (deleteaction.east) -- node[edu label] {Yes} (delete.west);
  \path [edu arrow] (edit) -- (submit);
  \path [edu arrow] (submit) -- (valid);
  \path [edu arrow] (valid) -- node[edu label] {Yes} (saveclass);
  \path [edu arrow] (valid.east) -- node[edu label] {No} (error.west);
  \path [edu arrow] (saveclass) -- (settings);
  \path [edu arrow] (settings) -- (refresh);
  \path [edu arrow] (delete) -- (removed);
  \path [edu arrow] (removed.south) |- (refresh.east);
  \path [edu arrow] (refresh) -- (end);
  \path [edu arrow] (error.south) |- (end.east);
\end{tikzpicture}%
}
